% ==========================================
% Copyright (c) 2026 CSharpKun & PositiveCHEEK
% This work is licensed under a Creative Commons Attribution-ShareAlike 4.0 International License.
% To view a copy of this license, visit http://creativecommons.org
% ==========================================

% Шаблон ВПР

\documentclass[vpr]{libre-oge}
\usepackage{lua-ul}
\usepackage{region-patterns}

\author{CSharpKun \& PositiveCHEEK}
\date{2026}
\setsubject{ПРЕДМЕТ}
\setklas{10}

\begin{document}

\begin{center}
	\vspace*{12mm}
	
	\textbf{Проверочная работа}
	
	\textbf{по \subject}
	
	\vspace{10mm}
	
	\textbf{\klas\space класс}
	
	\vspace{20mm}
	
	\textbf{Пояснение к проверочной работе}
	
	\vspace{10mm}
\end{center}

На выполнение работы по математике отводится два урока (не более 45 минут
каждый). Работа состоит из двух частей и включает в себя 17 заданий.
Обе части работы могут выполняться в один день с перерывом не менее 10 минут или
в разные дни.
При выполнении работы не разрешается пользоваться учебниками, рабочими
тетрадями, справочниками, калькулятором.
При необходимости можно пользоваться черновиком. Записи в черновике
проверяться и оцениваться не будут.

\newpage

\begin{center}
	\vspace*{10mm}
	\textbf{Инструкция по выполнению заданий части 1 проверочной работы}
	\vspace{10mm}
\end{center}

На выполнение заданий части 1 проверочной работы по математике отводится один
урок (не более 45 минут). Часть 1 включает в себя 12 заданий.

Ответы на задания запишите в поля ответов в тексте работы. Если Вы хотите
изменить ответ, зачеркните его и запишите рядом новый.

При выполнении работы не разрешается пользоваться учебниками, рабочими
тетрадями, справочниками, калькулятором.

При необходимости можно пользоваться черновиком. Записи в черновике
проверяться и оцениваться не будут.

Советуем выполнять задания в том порядке, в котором они даны. В целях экономии
времени пропускайте задание, которое не удаётся выполнить сразу, и переходите
к следующему. Если после выполнения работы у Вас останется время, то Вы сможете
вернуться к пропущенным заданиям.

\medskip

\begin{center}
	\textit{\textbf{Желаем успеха!}} 
\end{center}

\newpage

\begin{center}
\textbf{Часть 1}
\smallskip
\end{center}

\begin{task}{1}
Задание с кратким ответом
\vpranswerfield
\end{task}

\begin{task}{2}
Задание с кратким ответом
\vpranswerfield
\end{task}

\begin{task}{3}
Задание с кратким ответом
\vpranswerfield
\end{task}

\newpage

\begin{center}
	\vspace*{10mm}
	\textbf{Инструкция по выполнению заданий части 2 проверочной работы}
	\vspace{10mm}
\end{center}

На выполнение заданий части 2 проверочной работы по математике отводится один
урок (не более 45 минут). Часть 2 включает в себя 5 заданий.

В заданиях 13, 14, 16, 17 запишите решение и ответ в указанном месте. В задании 15
постройте график функции и ответьте на поставленный вопрос. Если Вы хотите изменить
ответ, зачеркните его и запишите рядом новый.

При выполнении работы не разрешается пользоваться учебниками, рабочими
тетрадями, справочниками, калькулятором.

При необходимости можно пользоваться черновиком. Записи в черновике
проверяться и оцениваться не будут.

Советуем выполнять задания в том порядке, в котором они даны. В целях экономии
времени пропускайте задание, которое не удаётся выполнить сразу, и переходите
к следующему. Если после выполнения работы у Вас останется время, то Вы сможете
вернуться к пропущенным заданиям.

\medskip

\begin{center}
	\textit{\textbf{Желаем успеха!}} 
\end{center}

\newpage

\begin{center}
	\textbf{Часть 2}
	\smallskip
\end{center}

\begin{task}{4}
Задание с развёрнутым ответом
\vpradvancedanswerfield
\end{task}

\begin{task}{5}
Задание с развёрнутым ответом
\vpradvancedanswerfield
\end{task}

\begin{task}{6}
Задание с развёрнутым ответом
\vpradvancedanswerfield
\end{task}

\end{document}
